The proposed descriptive superiority criterion is not met: CVT ties local spacing in completion count. Circle and both ellipse orientations fail spatial completion under all three methods. Both rectangle orientations complete under CVT and spacing, while the nearest-target control completes only the perpendicular rectangle. L reference and rotated L complete under CVT and spacing; nearest-target fails in both. All methods complete with the C opening toward the goal, whereas all three C-away tasks exhaust the original 150 s watchdog and fail. The C-away endpoint errors are 0.846, 0.496 and 0.818 m for CVT, nearest and spacing. These results show a task-orientation boundary under this common configuration, not a claim that C shapes are intrinsically impossible.

In total, 22 episodes settle and five reach the watchdog: ellipse-parallel/spacing, rotated-L/nearest and all three C-away cases. Every episode reaches geometric CONTACT\_READY and records physical contact and positive projected motion; none of those events suffices for completion. CVT and spacing achieve completion in the same five conditions but with different endpoint errors. These data support examining lateral motion and termination alongside deployment, rather than selecting a new feedback candidate solely from rank deficiency.
